\section{代码阅读法：先看数据如何流进 Trainer}
这节课的价值，在于把 Dataset 和 Trainer 两个核心抽象钉得很清楚。真正阅读框架代码时，最值得先跟的是：一个样本如何从 \texttt{get\_item} 进入 batch，再流入 Trainer 完成前向与反向传播。

\begin{knowledgebox}{阅读训练框架的最短路径}
\begin{itemize}
  \item 先看 Dataset 如何组织输入输出
  \item 再看 Trainer 如何消费 batch
  \item 最后才去看分布式策略如何包裹训练循环
\end{itemize}
\end{knowledgebox}

\subsection{本章小结}
理解训练框架，不一定要先啃最复杂的分布式代码；先看数据流和 Trainer 逻辑，往往更容易建立全局认识。

\newpage
